% =========================================================================
% Session-Tracked Agentic Software Development:
% From Spec-Driven Design to Unsupervised Development
%
% A framework for traceable, resumable, auditable LLM-driven software
% engineering via structured session management and semantic step tagging.
%
% Authors: NOAA EMC Global Workflow MCP Team
% Date: February 2026
% =========================================================================

\documentclass[11pt,letterpaper]{article}

% --- Core packages ---
\usepackage[utf8]{inputenc}
\usepackage[T1]{fontenc}
\usepackage[margin=1in]{geometry}
\usepackage{amsmath,amssymb}
\usepackage{graphicx}
\usepackage{booktabs}
\usepackage{enumitem}
\usepackage{microtype}

% --- Algorithm pseudocode ---
\usepackage{algorithm}
\usepackage{algorithmic}

% --- Code listings ---
\usepackage{listings}
\usepackage{xcolor}

% --- References ---
\usepackage[numbers,sort&compress]{natbib}

% --- Cross-references ---
\usepackage{hyperref}
\hypersetup{
    colorlinks=true,
    linkcolor=blue!70!black,
    citecolor=green!50!black,
    urlcolor=blue!70!black,
    pdftitle={Session-Tracked Agentic Software Development},
    pdfauthor={NOAA EMC Global Workflow MCP Team},
    pdfkeywords={Agentic AI, Spec-Driven Development, Session Tracking,
                 Unsupervised Development, LLM Software Engineering,
                 Model Context Protocol}
}

% --- Colored boxes for definitions ---
\usepackage{tcolorbox}
\tcbuselibrary{breakable}
\newtcolorbox{definition}[1][]{
    colback=blue!3,
    colframe=blue!50!black,
    fonttitle=\bfseries,
    title=#1,
    breakable
}
\newtcolorbox{keyinsight}[1][Key Insight]{
    colback=green!3,
    colframe=green!50!black,
    fonttitle=\bfseries,
    title=#1,
    breakable
}

% --- Code listing style ---
\lstdefinestyle{code}{
    basicstyle=\ttfamily\footnotesize,
    breaklines=true,
    frame=single,
    captionpos=b,
    numbers=left,
    numberstyle=\tiny\color{gray},
    keywordstyle=\color{blue!80!black}\bfseries,
    commentstyle=\color{green!50!black}\itshape,
    stringstyle=\color{red!70!black},
    backgroundcolor=\color{gray!3},
    showstringspaces=false
}
\lstset{style=code}

% --- Header/footer ---
\usepackage{fancyhdr}
\setlength{\headheight}{14pt}
\pagestyle{fancy}
\fancyhf{}
\lhead{\footnotesize Session-Tracked Agentic Software Development}
\rhead{\footnotesize NOAA EMC}
\cfoot{\thepage}

% =========================================================================
% TITLE
% =========================================================================
\title{%
    \textbf{Session-Tracked Agentic Software Development:} \\[0.3em]
    \textbf{From Spec-Driven Design to Unsupervised Development} \\[0.8em]
    \large A Framework for Traceable, Resumable, and Auditable \\
    \large LLM-Driven Software Engineering
}

\author{
    Terry McGuinness \\
    Environmental Modeling Center \\
    National Oceanic and Atmospheric Administration \\
    \texttt{Terry.McGuinness@noaa.gov}
}

\date{February 2026}

% =========================================================================
\begin{document}
% =========================================================================

\maketitle

% -----------------------------------------------------------------------
\begin{abstract}
Large language model (LLM) coding agents can generate, modify, and analyze
source code with increasing proficiency, yet their work remains largely
\emph{ephemeral}: each conversation starts from scratch, actions lack
structured audit trails, and there is no formal mechanism for an agent to
resume interrupted work across sessions or across different client
interfaces.  We present a \textbf{session-tracked agentic development
framework} that addresses these limitations through three interlocking
contributions:
(1)~\textbf{Spec-Driven Development (SDD)}, a methodology in which
machine-readable workflow specifications are the single source of truth
for what must be built, validated, and delivered;
(2)~\textbf{Session Management with semantic step tags}
(\texttt{research}, \texttt{design}, \texttt{implement},
\texttt{configure}, \texttt{validate}, \texttt{document},
\texttt{ingest}), which gives every agent action a durable, classified
audit record persisted to disk;
and (3)~a \textbf{two-layer intent architecture} that separates
\emph{protocol knowledge} (agent instruction files) from \emph{task
knowledge} (per-invocation plans), enabling a progression from
interactive supervised development (ISD) to fully unsupervised
development (USD).
We validate the framework through deployment on the NOAA Global Workflow
codebase, where a Phase~27 session completed 8/8 tracked steps---ingesting
9{,}155 new graph relationships across 383 shell-script nodes---in
5~minutes with a complete, query\-able audit trail.  The framework is
implemented atop the Model Context Protocol (MCP), making it portable
across VS~Code Copilot, Claude Desktop, Claude Code, and CI/CD pipelines
without modification.

\medskip\noindent
\textbf{Keywords:}
Agentic AI $\cdot$ Spec-Driven Development $\cdot$
Session Tracking $\cdot$ Unsupervised Development $\cdot$
LLM Software Engineering $\cdot$ Model Context Protocol $\cdot$
Knowledge Graphs $\cdot$ Retrieval-Augmented Generation
\end{abstract}

% -----------------------------------------------------------------------
\tableofcontents
\newpage

% =========================================================================
\section{Introduction}
\label{sec:intro}
% =========================================================================

\subsection{The Ephemeral Agent Problem}

LLM-based coding agents---GitHub Copilot, Claude Code, Cursor, Aider,
SWE-agent~\cite{yang2024swe}---have demonstrated remarkable capability on
benchmarks such as SWE-bench~\cite{jimenez2024swebench}.  Yet in practice,
deploying these agents on production codebases exposes a fundamental gap:
\emph{agent work is ephemeral}.

\begin{enumerate}[nosep]
    \item \textbf{No persistent state.}
          Each conversation is a blank slate.  If the user closes VS~Code,
          all context is lost; the agent cannot resume.
    \item \textbf{No structured audit trail.}
          Actions are buried in chat transcripts.  There is no
          machine-queryable record of \emph{what} was done, \emph{why},
          or \emph{in what order}.
    \item \textbf{No cross-client portability.}
          Work started in VS~Code cannot be continued in a terminal
          agent or a CI/CD pipeline without manual context transfer.
    \item \textbf{No intent classification.}
          An agent generating code and an agent researching architecture
          are indistinguishable in the conversation record.
\end{enumerate}

These limitations are acceptable for ad hoc question-answering but
untenable for \emph{systematic software development}---multi-step
engineering tasks that span hours or days, require traceable decisions,
and must survive interruptions.

\subsection{Our Approach: Session-Tracked Agentic Development}

We propose a framework that wraps LLM agent interactions in a
\textbf{session management layer} inspired by database transaction
semantics and workflow orchestration patterns~\cite{vanderaalst2004workflow}.
The core abstraction is:

\begin{definition}[Session]
A \textbf{session} is a bounded unit of development work that tracks a
named \emph{phase}, records individual \emph{steps} with semantic tags,
persists state to disk, and produces a durable \emph{history log} that
survives session completion.
\end{definition}

Sessions are managed by four MCP~\cite{anthropic2024mcp} tools:
\texttt{start\_sdd\_session}, \texttt{record\_sdd\_step},
\texttt{get\_sdd\_session} (resume), and \texttt{complete\_sdd\_session}.
These tools are exposed identically to every MCP-compatible client,
achieving cross-client portability by construction.

\subsection{Contributions}

\begin{enumerate}[nosep]
    \item \textbf{Spec-Driven Development (SDD)} methodology with
          machine-readable Markdown workflow specifications and phased
          execution (Section~\ref{sec:sdd}).
    \item \textbf{Phase~31 Session Model}: a lightweight, file-based
          session manager with seven semantic step tags and append-only
          history (Section~\ref{sec:session}).
    \item \textbf{Two-layer intent architecture} separating protocol
          instructions from task plans, enabling the progression from
          supervised to unsupervised execution
          (Section~\ref{sec:intent}).
    \item \textbf{Execution mode spectrum}: \texttt{dry\_run} $\to$
          \texttt{supervised} $\to$ \texttt{auto\_approved} $\to$
          \texttt{batch}, providing progressive trust levels
          (Section~\ref{sec:modes}).
    \item \textbf{Empirical validation} on the NOAA Global Workflow
          codebase with a complete 8-step session trace
          (Section~\ref{sec:evaluation}).
    \item \textbf{Multi-client MCP architecture} portable across
          VS~Code, Claude Desktop, terminal agents, and CI/CD
          (Section~\ref{sec:architecture}).
\end{enumerate}

\subsection{Paper Organization}

Section~\ref{sec:background} surveys related work.
Section~\ref{sec:sdd} presents the SDD methodology.
Section~\ref{sec:session} details the session management model.
Section~\ref{sec:intent} describes the two-layer intent architecture.
Section~\ref{sec:modes} defines the execution mode spectrum.
Section~\ref{sec:architecture} covers the MCP-based implementation.
Section~\ref{sec:evaluation} presents empirical results.
Section~\ref{sec:discussion} discusses limitations and future work.
Section~\ref{sec:conclusion} concludes.

% =========================================================================
\section{Background and Related Work}
\label{sec:background}
% =========================================================================

\subsection{LLM Coding Agents}

The rapid advancement of code-capable LLMs~\cite{chen2021codex,
roziere2024codellama} has spawned a generation of \emph{agentic} tools
that go beyond autocompletion to execute multi-step development tasks.
SWE-agent~\cite{yang2024swe} introduced an agent--computer interface
(ACI) design that achieves state-of-the-art performance on SWE-bench.
MetaGPT~\cite{hong2024metagpt} and ChatDev~\cite{qian2024chatdev} explore
multi-agent architectures where specialized roles (architect, developer,
tester) collaborate on software projects.

These systems demonstrate that LLMs can plan and execute complex
engineering tasks, but they share a common limitation:
\textbf{session state is in-memory and ephemeral}.  When the process
terminates, all execution context is lost.

\subsection{Tool Use and Reasoning}

ReAct~\cite{yao2023react} established the paradigm of interleaving
reasoning traces with tool invocations.
Toolformer~\cite{schick2023toolformer} showed that LLMs can learn to
call external tools autonomously.  Chain-of-thought
prompting~\cite{wei2022chain} demonstrated that explicit reasoning steps
improve task performance.

Our framework builds on these foundations but adds a critical missing
element: \textbf{tool calls that record their own execution in a
persistent, structured audit log}.  In our model, the tool
\texttt{record\_sdd\_step} is simultaneously an \emph{action} (the agent
records progress) and an \emph{observation} (the updated session state
is returned for context).

\subsection{Software Process Frameworks}

Agile methodologies~\cite{beck2001agile}, test-driven development
(TDD)~\cite{beck2003tdd}, and requirements engineering~\cite{sommerville2011requirements}
provide structured processes for human developers.
Continuous delivery~\cite{humble2010continuous} automates the build--test--deploy
pipeline.

SDD adapts these concepts for LLM agents: specifications replace user
stories, session state replaces sprint boards, and semantic step tags
replace story point categories.  The fundamental insight is that
\textbf{LLM agents need process discipline just as human teams do}---not
to constrain their generative capability, but to make their work
traceable, resumable, and auditable.

\subsection{Knowledge-Enhanced Code Intelligence}

Graph-based code analysis~\cite{allamanis2018survey} has a rich history.
GraphRAG~\cite{edge2024graphrag} extends retrieval-augmented
generation~\cite{lewis2020rag} with graph-derived community summaries.
Our earlier work demonstrated hybrid Neo4j+ChromaDB retrieval for
operational weather forecasting code, achieving $3.8\times$ improvement
in Precision@10 over text search baselines.

The session management framework presented here is \emph{orthogonal} to
RAG architecture: it tracks \emph{how} the agent uses these tools, not
\emph{what} the tools return.  However, the combination is powerful:
RAG provides context, and session tracking records how that context was
applied.

\subsection{Human-AI Interaction}

Guidelines for human--AI interaction~\cite{amershi2019guidelines}
emphasize that AI systems should support user control over system
actions and make clear what the system can do.
Studies of LLM coding tools~\cite{vaithilingam2022expectation} reveal
that developers value transparency in AI decision-making.

Our execution mode spectrum (Section~\ref{sec:modes}) directly implements
these principles: \texttt{supervised} mode provides per-step approval
gates, while \texttt{batch} mode delegates full autonomy only when trust
has been established through prior supervised sessions.

% =========================================================================
\section{Spec-Driven Development}
\label{sec:sdd}
% =========================================================================

\begin{definition}[SDD Principle]
\emph{``If it's not in the SDD, it doesn't get coded.''}
Every development task begins with a machine-readable specification that
defines objectives, steps, deliverables, and acceptance criteria before
any code is written.
\end{definition}

\subsection{Core Principles}

SDD is built on four foundational principles:

\begin{description}[nosep,leftmargin=1.5em]
    \item[P1: Specifications as First-Class Artifacts.]
    Development tasks are expressed as structured workflows in Markdown
    before implementation begins.  Specifications define \emph{what}
    (objectives), \emph{how} (validation criteria), \emph{when}
    (phase ordering), and \emph{who} (role assignment).

    \item[P2: Progressive Refinement.]
    Specifications evolve from high-level intent to detailed
    implementation through iterative passes, each adding precision
    while maintaining consistency with the parent specification.

    \item[P3: Health-Integrated Architecture.]
    All system components integrate health monitoring from initialization.
    The agent can verify preconditions (database connectivity, file
    system access) before attempting any step.

    \item[P4: Supervised Human--AI Collaboration.]
    AI agents propose implementations based on specifications; human
    experts review, refine, and approve.  The level of supervision is
    configurable along the execution mode spectrum.
\end{description}

\subsection{Workflow Specification Format}

SDD workflows are expressed as Markdown documents with structured
metadata.  We chose Markdown over YAML or JSON for three reasons:
(1)~human readability, (2)~compatibility with version control diffs,
and (3)~native rendering in documentation platforms.

A workflow specification contains:

\begin{lstlisting}[language={},caption={SDD Workflow Specification Structure},label=lst:spec]
# Phase Name
**Version**: X.Y  **Author**: Name  **Date**: YYYY-MM-DD

## Objective
High-level goal and scope definition.

## Steps
### Step 1: Step Name
**Type**: research | design | implement | configure | validate
**Deliverables**: Concrete outputs
**Acceptance Criteria**: Pass/fail conditions

### Step 2: ...
  (repeat for each step)

## Success Metrics
Quantifiable completion criteria.
\end{lstlisting}

\subsection{Phase Naming Convention}

Phases follow the pattern
\texttt{phase<N><letter>\_<descriptor>.md}, enabling hierarchical
decomposition:

\begin{itemize}[nosep]
    \item \texttt{phase27\_jjob\_script\_rag\_enhancement.md} --- parent phase
    \item \texttt{phase27a} through \texttt{phase27g} --- sub-phases A--G
    \item \texttt{phase27h\_multi\_collection\_routing.md} --- spin-off issue
\end{itemize}

The framework currently contains 31+ phases with sub-phases, spanning
infrastructure, tooling, knowledge ingestion, and methodology.

% =========================================================================
\section{Session Management Model}
\label{sec:session}
% =========================================================================

\subsection{Design Goals}

The session manager was designed with five goals:

\begin{enumerate}[nosep]
    \item \textbf{Lightweight}: No database required; file-based persistence.
    \item \textbf{Cross-conversation}: Survives server restarts and client reconnections.
    \item \textbf{Cross-client}: Works identically from VS~Code, Claude Desktop, or CLI.
    \item \textbf{Auditable}: Append-only history log for compliance and reconstruction.
    \item \textbf{Simple API}: Four tools covering the full lifecycle.
\end{enumerate}

\subsection{State Architecture}

Session state is maintained in two complementary files:

\begin{description}[nosep]
    \item[\texttt{active\_session.json}]
    Contains the current session state: session ID, phase name, start time,
    current step, completed steps array, skipped steps, and notes.
    Written on every \texttt{record\_sdd\_step} call.  Deleted on
    \texttt{complete\_sdd\_session}.

    \item[\texttt{history.jsonl}]
    Append-only event log in JSON Lines format.  Every lifecycle
    event---\texttt{started}, \texttt{step\_completed},
    \texttt{step\_skipped}, \texttt{resumed}, \texttt{completed},
    \texttt{abandoned}---is appended as an independent line.
    Never truncated or overwritten.
\end{description}

This dual-file architecture mirrors the write-ahead log (WAL) pattern
in database systems: \texttt{active\_session.json} is the ``current
state'' and \texttt{history.jsonl} is the ``redo log.''  Even if the
active session file is lost, the complete session can be reconstructed
from the history.

\subsection{Session Lifecycle}

Figure~\ref{fig:lifecycle} illustrates the session lifecycle as a state
machine.

\begin{figure}[h]
\centering
\begin{verbatim}
  start_sdd_session(phase, totalSteps, notes)
        |
        v
  +---------------+     record_sdd_step(step, name, tag, notes)
  |  in_progress  |<-----------------------------------------+
  +-------+-------+                                          |
          |                                                  |
          +--------------------------------------------------+
          |
          +---> complete_sdd_session(summary)
          |         |
          |         v
          |   +-----------+
          |   | completed |    (active_session.json deleted;
          |   +-----------+     history.jsonl retains all events)
          |
          +---> abandon_session(reason)
                    |
                    v
              +------------+
              | abandoned  |   (same file semantics as completed)
              +------------+
\end{verbatim}
\caption{Session lifecycle state machine.  The \texttt{get\_sdd\_session}
tool reads but does not modify state, enabling resume across conversations.}
\label{fig:lifecycle}
\end{figure}

\subsection{Semantic Step Tags}
\label{sec:tags}

Each recorded step carries a \textbf{semantic tag} drawn from a fixed
vocabulary of seven values:

\begin{table}[h]
\centering
\caption{Semantic step tags with intent classification and examples.}
\label{tab:tags}
\begin{tabular}{@{}llp{6cm}@{}}
\toprule
\textbf{Tag} & \textbf{Intent Class} & \textbf{Example} \\
\midrule
\texttt{research}   & Discovery    & Read existing code, explore architecture \\
\texttt{design}     & Planning     & Draft data model, choose algorithm \\
\texttt{implement}  & Construction & Write code, modify configuration \\
\texttt{configure}  & Setup        & Set environment variables, install dependencies \\
\texttt{validate}   & Verification & Run tests, check Neo4j counts, dry-run \\
\texttt{document}   & Knowledge    & Update README, write changelog entry \\
\texttt{ingest}     & Enrichment   & Run graph ingestion, update ChromaDB \\
\bottomrule
\end{tabular}
\end{table}

Tags serve three purposes:

\begin{enumerate}[nosep]
    \item \textbf{Intent classification.}  The tag tells \emph{what kind}
          of work the step represents, independent of its specific content.
    \item \textbf{Workflow analysis.}  Aggregating tags across sessions
          reveals development patterns (e.g., ratio of
          \texttt{implement} to \texttt{validate} steps).
    \item \textbf{Automation gating.}  In future work, the execution mode
          could auto-approve \texttt{research} and \texttt{validate}
          steps while requiring human approval for \texttt{implement}
          steps.
\end{enumerate}

\subsection{History Event Schema}

Each event appended to \texttt{history.jsonl} has the structure:

\begin{lstlisting}[language={},caption={History event for a step completion},label=lst:event]
{
  "sessionId": "session_2026-02-19_8aioyi",
  "phase": "phase27_jjob_script_rag_enhancement",
  "event": "step_completed",
  "step": 5,
  "name": "Run shell graph ingestion",
  "tag": "implement",
  "notes": "235 scripts processed, 383 ShellScript nodes,
            9155 new relationships",
  "timestamp": "2026-02-19T20:49:30.641Z"
}
\end{lstlisting}

The schema is deliberately flat---no nesting, no foreign keys---to
maximize grep-ability and ensure that every line is independently
parseable.  This design choice favors operational simplicity over
relational normalization.

\subsection{Cross-Conversation Resume}

The \texttt{get\_sdd\_session} tool reads \texttt{active\_session.json}
from disk and returns the full session state, including all completed
steps.  Because the file persists across server restarts, an agent in a
\emph{new} conversation can resume precisely where the previous
conversation left off:

\begin{algorithm}
\caption{Cross-conversation resume protocol}
\label{alg:resume}
\begin{algorithmic}[1]
\STATE Agent starts in new conversation
\STATE Agent calls \texttt{get\_sdd\_session()}
\IF{active session exists}
    \STATE Read completed steps $\{s_1, \ldots, s_k\}$
    \STATE Read total steps $N$ from specification
    \STATE Continue from step $k+1$
\ELSE
    \STATE No active session; start fresh or report idle
\ENDIF
\end{algorithmic}
\end{algorithm}

This mechanism requires \textbf{no conversation memory, no external
database, and no user intervention}.  The agent discovers its own
state through a tool call.

% =========================================================================
\section{Two-Layer Intent Architecture}
\label{sec:intent}
% =========================================================================

A critical challenge in agentic development is communicating \emph{both}
the development methodology (``how we build software'') and the specific
task (``what to build right now'') to the LLM.  We address this with a
\textbf{two-layer intent stack}:

\begin{table}[h]
\centering
\caption{Two-layer intent architecture.}
\label{tab:intent}
\begin{tabular}{@{}lllp{4.5cm}@{}}
\toprule
\textbf{Layer} & \textbf{Scope} & \textbf{Mechanism} & \textbf{Content} \\
\midrule
Protocol & Persistent & Instruction files & SDD methodology, tool usage,
    session lifecycle, conventions \\
Task     & Per-invocation & \texttt{/plan} or spec & Phase objectives,
    step details, acceptance criteria \\
\bottomrule
\end{tabular}
\end{table}

\subsection{Protocol Layer: Instruction Files}

Instruction files (e.g., \texttt{.github/instructions/*.md}) are loaded
into the LLM's system context on every invocation.  They encode:

\begin{itemize}[nosep]
    \item \textbf{SDD methodology}: ``If it's not in the SDD, it doesn't
          get coded.''
    \item \textbf{Session protocol}: Start sessions before work; record
          steps with semantic tags; complete sessions with summaries.
    \item \textbf{Tool usage}: Parameter names, required vs.\ optional
          arguments, common workflows.
    \item \textbf{Conventions}: ASCII-only console output, 2-space
          indentation, SPOT (Single Point of Truth) for configuration.
\end{itemize}

The LLM internalizes these instructions as \emph{behavioral constraints}
that apply to every task, regardless of specifics.  This is analogous
to a developer internalizing a team's coding standards.

\subsection{Task Layer: Plans and Specifications}

The task layer provides per-invocation context through two mechanisms:

\begin{description}[nosep]
    \item[Explicit plans (\texttt{/plan}):]
    The user provides a natural-language plan that the LLM decomposes
    into an internal task list.  The plan is ephemeral---it exists only
    in the current conversation context.

    \item[Workflow specifications:]
    A Markdown specification (Section~\ref{sec:sdd}) provides structured
    step definitions.  The agent reads the spec via
    \texttt{get\_sdd\_workflow} and auto-generates its own internal plan
    from the structured content.
\end{description}

\begin{keyinsight}
The protocol layer teaches the agent \emph{how to work}; the task layer
teaches the agent \emph{what to build}.  Together, they provide
sufficient intent for the agent to self-plan, self-tag, self-validate,
and self-resume without human prompting at each step.
\end{keyinsight}

\subsection{Intent Propagation}

The complete intent stack for any given agent action is:

\begin{equation}
\label{eq:intent}
\text{Intent}(a) = \underbrace{\text{Spec}(p)}_{\text{What to build}}
    \oplus \underbrace{\text{Tag}(s)}_{\text{What kind of work}}
    \oplus \underbrace{\text{Name}(s) + \text{Notes}(s)}_{\text{Specifics}}
\end{equation}

\noindent where $p$ is the phase specification, $s$ is the current step,
and $\oplus$ denotes context concatenation.  Each component adds a
layer of semantic precision:

\begin{itemize}[nosep]
    \item \textbf{Spec}: ``Phase 27: J-Job Script RAG Enhancement''
    \item \textbf{Tag}: \texttt{implement} (this is a construction step)
    \item \textbf{Name}: ``Run shell graph ingestion''
    \item \textbf{Notes}: ``235 scripts processed, 383 ShellScript nodes,
          9{,}155 new relationships''
\end{itemize}

This layered structure means that even if the name and notes are terse,
the combination of spec~+~tag~+~name provides enough context to
reconstruct the full intent post hoc.

% =========================================================================
\section{Execution Mode Spectrum}
\label{sec:modes}
% =========================================================================

The framework defines four execution modes with \textbf{progressive
trust levels}:

\begin{table}[h]
\centering
\caption{Execution mode spectrum.  Trust increases left to right.}
\label{tab:modes}
\begin{tabular}{@{}lcccl@{}}
\toprule
\textbf{Mode} & \textbf{Side Effects} & \textbf{Human} & \textbf{Audit} & \textbf{Use Case} \\
\midrule
\texttt{dry\_run}       & None     & None       & Preview only  & Planning, validation \\
\texttt{supervised}     & Yes      & Per-step   & Full          & Active development \\
\texttt{auto\_approved} & Yes      & Manifest   & Full          & Trusted workflows \\
\texttt{batch}          & Yes      & None       & Full          & CI/CD pipelines \\
\bottomrule
\end{tabular}
\end{table}

\subsection{Step Classification}

Steps are classified by their side-effect potential:

\begin{description}[nosep]
    \item[Read-only] (\texttt{research}, \texttt{validate}):
    Auto-execute in all modes.  These steps query databases, read files,
    and check conditions without modifying system state.

    \item[Side-effect] (\texttt{implement}, \texttt{configure},
    \texttt{ingest}):
    Require approval in \texttt{supervised} mode, pre-approved by
    manifest in \texttt{auto\_approved} mode, or fully autonomous in
    \texttt{batch} mode.

    \item[Documentation] (\texttt{document}, \texttt{design}):
    Typically auto-approved since they produce artifacts (Markdown,
    diagrams) that are easily reviewed and reverted.
\end{description}

\subsection{From ISD to USD}

The execution mode spectrum enables a natural progression from
\textbf{Interactive Supervised Development (ISD)} to
\textbf{Unsupervised Development (USD)}:

\begin{enumerate}[nosep]
    \item \textbf{ISD Phase}: The developer runs sessions in
          \texttt{supervised} mode, reviewing each side-effect step.
          The session history accumulates a record of what actions were
          approved and which patterns are trustworthy.
    \item \textbf{Trust Calibration}: After multiple successful supervised
          sessions, the developer identifies step types and workflow
          patterns that consistently produce correct results.
    \item \textbf{USD Phase}: Trusted patterns are codified in an
          auto-approval manifest.  The agent executes autonomously
          within these boundaries while still recording every step to
          the audit log.
\end{enumerate}

\begin{keyinsight}
USD is not ``unsupervised'' in the sense of unmonitored.
Every step is still recorded with full semantic tagging and notes.
The difference is \emph{when} the human reviews: in ISD, review is
synchronous (before each step); in USD, review is asynchronous
(after completion, via the audit log).
\end{keyinsight}

% =========================================================================
\section{Implementation Architecture}
\label{sec:architecture}
% =========================================================================

\subsection{Technology Stack}

The framework is implemented as part of a unified MCP server:

\begin{table}[h]
\centering
\caption{Technology stack.}
\label{tab:stack}
\begin{tabular}{@{}lll@{}}
\toprule
\textbf{Component} & \textbf{Technology} & \textbf{Purpose} \\
\midrule
MCP Server       & Node.js (ES Modules) & Tool orchestration via MCP SDK \\
Session Manager  & \texttt{SessionManager.js} & Session lifecycle and history \\
Graph Database   & Neo4j 5.x & Code structure relationships \\
Vector Database  & ChromaDB 0.4+ & Semantic search embeddings \\
Transport        & stdio / Streamable HTTP & MCP client communication \\
State Storage    & JSON + JSONL files & Session persistence \\
\bottomrule
\end{tabular}
\end{table}

\subsection{SessionManager Implementation}

The \texttt{SessionManager} class (260~LOC) implements the full session
lifecycle with seven public methods:

\begin{lstlisting}[language=Java,caption={SessionManager API},label=lst:api]
class SessionManager {
  startSession(phaseName, options)   // -> session state
  recordStep(stepNumber, name, tag, notes) // -> updated state
  skipStep(stepNumber, reason)       // -> updated state
  getSessionState()                  // -> state | null
  resumeSession()                    // -> state (with timestamp)
  completeSession(summary)           // -> final state
  abandonSession(reason)             // -> final state
}
\end{lstlisting}

Key implementation details:

\begin{itemize}[nosep]
    \item \textbf{Single active session constraint}.
          \texttt{startSession} throws if an active session already
          exists.  This prevents interleaved sessions which would
          complicate the audit trail.
    \item \textbf{Duplicate detection.}
          \texttt{recordStep} rejects duplicate step numbers, preventing
          accidental re-recording.
    \item \textbf{Tag validation.}
          Unknown tags produce a warning but do not block recording,
          allowing graceful evolution of the tag vocabulary.
    \item \textbf{Auto-parsing step count.}
          \texttt{startSession} reads the workflow specification file
          and counts \texttt{\#\#\# Step N:} patterns to auto-populate
          \texttt{totalSteps}.
    \item \textbf{Duration calculation.}
          \texttt{completeSession} computes wall-clock duration from
          start to completion timestamps.
\end{itemize}

\subsection{MCP Tool Registration}

The four session tools are registered with the MCP server alongside
38+ other tools (code analysis, semantic search, compliance checking):

\begin{table}[h]
\centering
\caption{SDD session tools: MCP registration.}
\label{tab:tools}
\begin{tabular}{@{}llp{5.5cm}@{}}
\toprule
\textbf{Tool} & \textbf{Required Params} & \textbf{Description} \\
\midrule
\texttt{start\_sdd\_session}    & \texttt{phase} &
    Start tracked session; optional \texttt{totalSteps}, \texttt{notes} \\
\texttt{record\_sdd\_step}      & \texttt{step}, \texttt{name} &
    Record completion; optional \texttt{tag}, \texttt{notes} \\
\texttt{get\_sdd\_session}      & (none) &
    Return active session state for resume \\
\texttt{complete\_sdd\_session} & (none) &
    Complete session; optional \texttt{summary} \\
\bottomrule
\end{tabular}
\end{table}

\subsection{Multi-Client Portability}

Because session state is persisted to disk (not in MCP server memory),
the framework is inherently multi-client:

\begin{enumerate}[nosep]
    \item A developer starts a session in \textbf{VS~Code Copilot}
          (stdio transport) and completes steps~1--3.
    \item The developer closes VS~Code.
    \item Later, a \textbf{Claude Code} terminal agent calls
          \texttt{get\_sdd\_session()}, discovers steps~1--3 complete,
          and continues from step~4.
    \item A \textbf{CI/CD pipeline} calls
          \texttt{complete\_sdd\_session()} after automated validation.
\end{enumerate}

This is possible because all clients connect to the same MCP server
instance (or containers reading the same volume mount), and the session
state is a simple JSON file on a shared filesystem.

\subsection{Docker MCP Gateway}

For production deployment, the MCP server runs inside a Docker container
behind the Docker MCP Gateway (Streamable HTTP on port~18888).  The
\texttt{sdd\_framework/} directory is volume-mounted, ensuring session
state survives container restarts:

\begin{lstlisting}[language=bash,caption={Volume mount for session persistence}]
volumes:
  - /path/to/sdd_framework:/app/sdd_framework   # state persists
  - /path/to/supported_repos:/app/supported_repos # read-only
\end{lstlisting}

% =========================================================================
\section{Empirical Evaluation}
\label{sec:evaluation}
% =========================================================================

\subsection{Case Study: Phase 27F-G Shell Graph Ingestion}

We evaluate the framework through a complete session trace from
Phase~27 of the NOAA Global Workflow RAG enhancement project.  This
phase involved ingesting shell script structure (J-Jobs, ex-scripts,
environment variables) into the Neo4j graph database.

\subsubsection{Session Metadata}

\begin{table}[h]
\centering
\caption{Phase 27F-G session metadata.}
\label{tab:session}
\begin{tabular}{@{}ll@{}}
\toprule
\textbf{Attribute} & \textbf{Value} \\
\midrule
Session ID        & \texttt{session\_2026-02-19\_8aioyi} \\
Phase             & \texttt{phase27\_jjob\_script\_rag\_enhancement} \\
Total Steps       & 8 (auto-parsed from specification) \\
Completed Steps   & 8/8 (100\%) \\
Duration          & 5 minutes \\
Client            & VS~Code Copilot (stdio) \\
\bottomrule
\end{tabular}
\end{table}

\subsubsection{Step-by-Step Trace}

Table~\ref{tab:steps} shows the complete step trace as recorded in
\texttt{history.jsonl}.

\begin{table}[h]
\centering
\caption{Complete step trace for Phase 27F-G.}
\label{tab:steps}
\small
\begin{tabular}{@{}cllp{5.8cm}@{}}
\toprule
\textbf{\#} & \textbf{Tag} & \textbf{Name} & \textbf{Key Metrics} \\
\midrule
1 & \texttt{implement}  & Fix shell graph password  & SPOT-compliant default \\
2 & \texttt{implement}  & Add CLI flags & \texttt{--dry-run}, \texttt{--clear}, \texttt{--verbose} \\
3 & \texttt{implement}  & Fix bridge paths & Duplicate session, path refactor \\
4 & \texttt{validate}   & Dry-run shell graph & 235 scripts, 0 errors \\
5 & \texttt{implement}  & Run ingestion & 383 nodes, 9{,}155 relationships \\
6 & \texttt{implement}  & Re-run bridges & 8 cross-language edges \\
7 & \texttt{validate}   & Validate Neo4j counts & 40{,}413 total nodes (+206) \\
8 & \texttt{validate}   & Validate MCP tools & 5/6 tests pass, 1 miss \\
\bottomrule
\end{tabular}
\end{table}

\subsubsection{Tag Distribution}

The session's tag distribution reveals the development pattern:

\begin{itemize}[nosep]
    \item \texttt{implement}: 5 steps (62.5\%)---majority of effort in code changes
    \item \texttt{validate}: 3 steps (37.5\%)---significant validation effort
    \item \texttt{research}: 0 steps---the phase spec provided sufficient context
\end{itemize}

This ratio (roughly 2:1 implement-to-validate) is characteristic of a
well-specified phase where the agent does not need to ``explore''---the
specification told it exactly what to build, and the validation steps
confirmed correctness.

\subsubsection{Quantitative Results}

\begin{table}[h]
\centering
\caption{Knowledge graph metrics before and after Phase 27F-G.}
\label{tab:metrics}
\begin{tabular}{@{}lrrl@{}}
\toprule
\textbf{Metric} & \textbf{Before} & \textbf{After} & \textbf{Delta} \\
\midrule
Total nodes              & 40{,}207 & 40{,}413 & +206 \\
ShellScript nodes        & 0        & 383      & +383 \\
ShellFunction nodes      & 0        & 63       & +63  \\
EnvironmentVariable nodes& n/a      & 2{,}489  & +2{,}489 \\
SOURCES relationships     & 148      & 393      & +245 \\
INVOKES relationships     & 4        & 352      & +348 \\
DEPENDS\_ON\_ENV          & 6{,}007  & 7{,}225  & +1{,}218 \\
DEFINES relationships     & 9{,}690  & 9{,}753  & +63  \\
Total new relationships  & ---      & ---      & +9{,}155 \\
\bottomrule
\end{tabular}
\end{table}

\subsubsection{Audit Trail Verification}

The complete session is reconstructable from \texttt{history.jsonl}
alone.  We verified this by:

\begin{enumerate}[nosep]
    \item Parsing all events with \texttt{sessionId =
          session\_2026-02-19\_8aioyi}
    \item Confirming \texttt{started} event with correct phase name
    \item Confirming 8 \texttt{step\_completed} events in order
    \item Confirming \texttt{completed} event with duration and step counts
    \item Total: 10 events (1 started + 8 step\_completed + 1 completed)
\end{enumerate}

\subsection{Session Management Validation}

Prior to the Phase~27 session, two validation sessions tested the
session lifecycle:

\begin{enumerate}[nosep]
    \item \textbf{Phase~31 review test}
          (\texttt{session\_2026-02-18\_8p8fsr}): 2 steps completed,
          1 step skipped, verifying the skip and completion mechanisms.
    \item \textbf{Health check test}
          (\texttt{session\_2026-02-19\_x8kyst}): 2 steps completed
          with an explicit \texttt{resumed} event, verifying the
          cross-conversation resume capability.
\end{enumerate}

All three sessions are permanently recorded in \texttt{history.jsonl}
(17 total events).

\subsection{Comparison with Untracked Development}

Without session tracking, the same Phase~27 work would have produced:

\begin{itemize}[nosep]
    \item A chat transcript (ephemeral, not machine-queryable)
    \item Git commits (capture \emph{what} changed, not \emph{why} or
          \emph{in what order})
    \item No structured record of the dry-run step that validated
          correctness before the live ingestion
    \item No tag-based analysis of implementation vs.\ validation effort
\end{itemize}

The session trace adds approximately 3~KB of structured metadata
(10 JSONL events) that makes the entire development process
reconstructable and analyzable.

% =========================================================================
\section{Discussion}
\label{sec:discussion}
% =========================================================================

\subsection{Limitations}

\subsubsection{Single Active Session}

The current design enforces one active session at a time.  This prevents
interleaved work (e.g., pausing Phase~27 to fix an urgent bug in a
different phase).  The workaround is to \texttt{abandon} the current
session and start a new one, but this loses the in-progress state.

\textbf{Proposed solution}: Named session stacks, where sessions can be
suspended and resumed by name, similar to Git stash.

\subsubsection{Completion Snapshot Loss}

When \texttt{complete\_sdd\_session} is called, the active session file
is deleted.  While all events persist in \texttt{history.jsonl}, the
assembled session object (with its complete step array) is lost.  To
reconstruct the full session state, one must filter and reassemble events
from the history file.

\textbf{Proposed solution}: Write a snapshot entry to
\texttt{history.jsonl} containing the full session state before deletion.

\subsubsection{No Step Reordering}

Steps must be recorded in increasing numerical order (enforced by
duplicate detection).  Real-world development sometimes requires
revisiting earlier steps.  The current design handles this by recording
a new step with a higher number and explanatory notes.

\subsubsection{Tag Vocabulary Evolution}

The fixed tag vocabulary (\texttt{research}, \texttt{design},
\texttt{implement}, \texttt{configure}, \texttt{validate},
\texttt{document}, \texttt{ingest}) may not capture all development
activities.  Unknown tags produce warnings but are accepted, providing
a graceful evolution path.

\subsection{Generalizability}

While developed for NOAA Global Workflow, the framework is
domain-agnostic.  Any codebase that can be described by Markdown
specifications can use SDD session tracking.  The MCP transport layer
means any LLM with MCP support can participate.

The semantic step tags are also generalizable: \texttt{research},
\texttt{design}, \texttt{implement}, \texttt{validate}, and
\texttt{document} apply to virtually any software project.  The
\texttt{configure} and \texttt{ingest} tags are more specific to
infrastructure and knowledge-base maintenance but could be replaced with
domain-appropriate tags.

\subsection{Comparison with Existing Approaches}

\begin{table}[h]
\centering
\caption{Comparison with existing development frameworks.}
\label{tab:comparison}
\small
\begin{tabular}{@{}lccccc@{}}
\toprule
\textbf{Feature} &
\rotatebox{70}{\textbf{SDD (this work)}} &
\rotatebox{70}{\textbf{SWE-agent}} &
\rotatebox{70}{\textbf{MetaGPT}} &
\rotatebox{70}{\textbf{Claude Code}} &
\rotatebox{70}{\textbf{GitHub Actions}} \\
\midrule
Persistent sessions       & \checkmark &            &            &            & \\
Semantic step tags        & \checkmark &            &            &            & \\
Cross-client resume       & \checkmark &            &            &            & \\
Append-only audit log     & \checkmark &            &            &            & \checkmark \\
Progressive trust modes   & \checkmark &            &            & \checkmark & \\
Spec-first methodology    & \checkmark &            & \checkmark &            & \checkmark \\
Multi-agent support       &            &            & \checkmark &            & \\
Benchmark evaluation      &            & \checkmark & \checkmark &            & \\
\bottomrule
\end{tabular}
\end{table}

SWE-agent~\cite{yang2024swe} and MetaGPT~\cite{hong2024metagpt} excel at
benchmark tasks but lack persistent session management.
Claude Code~\cite{anthropic2025claudecode} provides progressive trust
(auto-accept patterns) but does not persist state across invocations.
GitHub Actions provides audit trails and spec-driven workflows but
operates at a different abstraction level (CI/CD pipelines rather than
interactive development).

Our framework uniquely combines persistent sessions, semantic tagging,
cross-client portability, and progressive trust in a single lightweight
architecture.

\subsection{Future Work}

\begin{description}[nosep]
    \item[Multi-session support:]
    Named session stacks for concurrent work across multiple phases.

    \item[Completion snapshots:]
    Write full session state to history on completion, eliminating
    reconstruction overhead.

    \item[Tag-based auto-approval:]
    In \texttt{auto\_approved} mode, automatically approve steps tagged
    \texttt{research} or \texttt{validate} while gating
    \texttt{implement} steps.

    \item[Session analytics dashboard:]
    Aggregate tag distributions, duration patterns, and completion rates
    across all historical sessions.

    \item[Multi-agent coordination:]
    Multiple agents working on the same phase from different clients,
    with conflict detection based on step-number reservation.

    \item[Stale session detection:]
    Automatically flag sessions with no activity beyond a configurable
    TTL, enabling cleanup of abandoned sessions.

    \item[LLM self-evaluation:]
    Use the session audit trail as fine-tuning data: train models to
    predict which tag and notes format produces the most useful history.
\end{description}

% =========================================================================
\section{Conclusion}
\label{sec:conclusion}
% =========================================================================

We have presented a session-tracked agentic development framework that
transforms ephemeral LLM coding interactions into structured, resumable,
and auditable engineering processes.  The framework's three core
contributions---Spec-Driven Development methodology, the Phase~31
session management model with semantic step tags, and the two-layer
intent architecture---address fundamental gaps in current LLM agent
tooling.

Empirical evaluation on the NOAA Global Workflow codebase demonstrates
the framework's practical value: a Phase~27 session completed 8/8
tracked steps in 5~minutes, ingesting 9{,}155 new knowledge graph
relationships across 383 shell-script nodes, with a complete
machine-queryable audit trail of 10~events in the append-only history
log.

The execution mode spectrum (\texttt{dry\_run} $\to$ \texttt{supervised}
$\to$ \texttt{auto\_approved} $\to$ \texttt{batch}) provides a
principled path from interactive supervised development to unsupervised
autonomy, with trust calibrated through experience rather than
configuration.

Perhaps most importantly, the entire framework is implemented as four
MCP tools totaling 260~lines of code, persisting state in two files on
a shared filesystem.  This simplicity is a feature, not a limitation:
it means the framework can be adopted incrementally, runs on any platform
with a filesystem, and requires no external services beyond what the LLM
agent already uses.

\textbf{Key insight}: The gap between current LLM agent capability and
production-grade software engineering is not primarily a capability
gap---it is a \emph{process} gap.  LLMs can write code, analyze
architectures, and validate results.  What they lack is a structured
process for recording \emph{what they did, why they did it, and where
they left off}.  Session-tracked agentic development fills this gap.

% =========================================================================
\section*{Acknowledgments}
% =========================================================================

This work was developed at the NOAA Environmental Modeling Center as part
of the Global Workflow modernization effort.  The authors thank the
operational forecasting staff for domain expertise and the Parallel Works
team for HPC infrastructure support.

% =========================================================================
% REFERENCES
% =========================================================================
\bibliographystyle{plainnat}
\bibliography{references}

% =========================================================================
\newpage
\appendix
% =========================================================================

\section{Complete SessionManager.js API Reference}
\label{app:api}

\begin{table}[h]
\centering
\caption{SessionManager method signatures and semantics.}
\label{tab:api}
\small
\begin{tabular}{@{}lp{4cm}p{5cm}@{}}
\toprule
\textbf{Method} & \textbf{Parameters} & \textbf{Side Effects} \\
\midrule
\texttt{startSession} &
    \texttt{phaseName}: string \newline
    \texttt{options}: \{totalSteps, notes\} &
    Writes \texttt{active\_session.json} \newline
    Appends \texttt{started} to history \\
\midrule
\texttt{recordStep} &
    \texttt{stepNumber}: int \newline
    \texttt{name}: string \newline
    \texttt{tag}: string \newline
    \texttt{notes}: string &
    Updates \texttt{active\_session.json} \newline
    Appends \texttt{step\_completed} to history \\
\midrule
\texttt{skipStep} &
    \texttt{stepNumber}: int \newline
    \texttt{reason}: string &
    Updates \texttt{active\_session.json} \newline
    Appends \texttt{step\_skipped} to history \\
\midrule
\texttt{getSessionState} &
    (none) &
    Read-only; returns session or null \\
\midrule
\texttt{resumeSession} &
    (none) &
    Updates \texttt{lastActivityAt} \newline
    Appends \texttt{resumed} to history \\
\midrule
\texttt{completeSession} &
    \texttt{summary}: string &
    Appends \texttt{completed} to history \newline
    \textbf{Deletes} \texttt{active\_session.json} \\
\midrule
\texttt{abandonSession} &
    \texttt{reason}: string &
    Appends \texttt{abandoned} to history \newline
    \textbf{Deletes} \texttt{active\_session.json} \\
\bottomrule
\end{tabular}
\end{table}

\section{History Event Types}
\label{app:events}

\begin{table}[h]
\centering
\caption{All event types recorded in \texttt{history.jsonl}.}
\label{tab:events}
\begin{tabular}{@{}lp{7cm}@{}}
\toprule
\textbf{Event} & \textbf{Additional Fields} \\
\midrule
\texttt{started}         & phase \\
\texttt{step\_completed}  & step, name, tag, notes \\
\texttt{step\_skipped}    & step, reason \\
\texttt{resumed}          & (none beyond base fields) \\
\texttt{completed}        & summary, completedSteps, skippedSteps,
                            totalSteps, duration \\
\texttt{abandoned}        & reason, completedSteps \\
\bottomrule
\end{tabular}
\end{table}

All events share base fields: \texttt{sessionId}, \texttt{phase},
\texttt{event}, \texttt{timestamp}.

\section{Semantic Tag Vocabulary}
\label{app:tags}

The tag vocabulary is defined as a constant in \texttt{SessionManager.js}:

\begin{lstlisting}[language=Java]
const VALID_TAGS = [
    'research',    // Discovery: read code, explore architecture
    'design',      // Planning: draft models, choose algorithms
    'implement',   // Construction: write code, modify config
    'configure',   // Setup: environment, dependencies, infra
    'validate',    // Verification: tests, counts, dry-runs
    'document',    // Knowledge: README, changelog, comments
    'ingest'       // Enrichment: graph ingestion, vector updates
];
\end{lstlisting}

Tags are validated with a warning (not an error) for unknown values,
allowing the vocabulary to evolve without breaking existing sessions.

\section{Phase 27F-G Complete History Trace}
\label{app:trace}

The following is the verbatim content of \texttt{history.jsonl} entries
for session \texttt{session\_2026-02-19\_8aioyi}, demonstrating the
complete audit trail:

\bigskip
\noindent\textbf{Event 1} (session start):
\begin{lstlisting}[language={}]
{"sessionId":"session_2026-02-19_8aioyi",
 "phase":"phase27_jjob_script_rag_enhancement",
 "event":"started",
 "timestamp":"2026-02-19T20:46:46.502Z"}
\end{lstlisting}

\noindent\textbf{Event 2} (step 1 --- implement):
\begin{lstlisting}[language={}]
{"sessionId":"session_2026-02-19_8aioyi",
 "phase":"phase27_jjob_script_rag_enhancement",
 "event":"step_completed", "step":1,
 "name":"Fix shell graph password default",
 "tag":"implement",
 "notes":"Changed NEO4J_PASSWORD default from 'password'
   to 'gfsworkflow2025'. SPOT compliance.",
 "timestamp":"2026-02-19T20:47:15.678Z"}
\end{lstlisting}

\noindent\textbf{Event 6} (step 5 --- implement, key ingestion):
\begin{lstlisting}[language={}]
{"sessionId":"session_2026-02-19_8aioyi",
 "phase":"phase27_jjob_script_rag_enhancement",
 "event":"step_completed", "step":5,
 "name":"Run shell graph ingestion",
 "tag":"implement",
 "notes":"235 scripts processed (0 errors), 383 ShellScript
   nodes (89 J-Jobs, 130 ex-scripts), 2489 EnvironmentVariables,
   63 ShellFunctions, 1 ConfigFile. Relationships: 393 SOURCES,
   352 INVOKES, 1184 EXPORTS, 7225 DEPENDS_ON_ENV,
   1 READS_CONFIG = 9,155 total new relationships.",
 "timestamp":"2026-02-19T20:49:30.641Z"}
\end{lstlisting}

\noindent\textbf{Event 10} (session completion):
\begin{lstlisting}[language={}]
{"sessionId":"session_2026-02-19_8aioyi",
 "phase":"phase27_jjob_script_rag_enhancement",
 "event":"completed",
 "completedSteps":8, "skippedSteps":0, "totalSteps":8,
 "duration":"5m",
 "timestamp":"2026-02-19T20:52:02.928Z"}
\end{lstlisting}

% =========================================================================
\end{document}
